\section{Related Work}
\begin{table*}[h]
\centering
\small
\setlength{\tabcolsep}{3pt}
\begin{tabular}{lccc}
\toprule
\textbf{Frameworks} & \textbf{Training-Rollout Decoupled?} & \textbf{Rootless Sandbox?} & \textbf{Scaffold-Independent?} \\
\midrule
SkyRL-Agent~\citep{cao2025skyrlagent} & \xmark & \xmark & \checkmark \\
VeRL-Tool~\citep{jiang2025verltool} & \xmark & \xmark & \checkmark \\
Agent Lightning~\citep{luo2025agentlightning} & \xmark & \xmark & \xmark \\
rLLM~\citep{tan2025rllm} & \xmark & \xmark & \checkmark \\
GEM~\citep{liu2025gem} & \xmark & \xmark & \checkmark \\
\midrule
\rowcolor{gray!30} \textbf{\prorlagent (Ours)} & \checkmark & \checkmark & \checkmark \\
\bottomrule
\end{tabular}
\caption{\textbf{Comparison of \prorlagent with existing frameworks for multi-turn agent RL.} \prorlagent decouples rollout from training, supports rootless sandboxing for shared HPC environments, and is independent of any specific training framework.}
\label{tab:framework_comparison}
\end{table*}

\textbf{Multi-turn RL for LLM Agents.}
Reinforcement learning has been highly effective for improving single-turn reasoning such as mathematics, logic, and coding~\citep{shao2024deepseekmath, guo2025deepseek, hu2025openreasonerzero, zhang2026nemotronresearchtooln}. 
Building on this progress, recent work has extended RL to multi-turn agentic settings, where agents interact with external environments over long horizons~\citep{cao2025skyrl, luo2025deepswe, gao2025beyond, li2025torl, jin2025searchr1,wang2025vagen,ragenv2026collapse}. 
In these settings, a multi-turn agent is naturally formulated as a POMDP~\citep{kaelbling1998pomdp}, where agent produces actions through tool calls~\citep{yao2022react, wang2024codeact,patil2025the,zhang2024offline} and receives environment observations at each step.
As tasks become more complex, multi-turn rollouts often span dozens of steps in diverse environments, such as code repositories~\citep{jimenez2024swebench, jain2025r2egym}, web browsers~\citep{zhou2023webarena}, and even computer operating systems~\citep{xie2024osworld}.

As a result, the infrastructure required to generate, manage, and evaluate these rollouts at scale has become a major bottleneck for RL training.
This bottleneck slows both training and the deployment of RL agents. \prorlagent is designed to address this challenge by decoupling the full lifecycle of multi-turn agent rollout from the training stack, allowing researchers and practitioners to focus on training algorithms and agent design.

\textbf{Agent RL Infrastructures.}
A growing body of work has begun to address the challenges of scalable RL training for agents, including support for diverse tool integration~\citep{jiang2025verltool, li2025torl}, flexible environment abstractions~\citep{liu2025gem, tan2025rllm}, and efficient rollout scheduling~\citep{cao2025skyrlagent}.
Yet across these frameworks, rollout orchestration, including environment lifecycle management, tool execution, trajectory collection, and evaluation, remains implemented as an in-process library within the training loop. 
Under this design, adopting a new training backend often requires re-implementing or porting the entire rollout stack. This tight coupling makes rollout infrastructure a major source of friction in multi-turn agent RL, often demanding more engineering effort than the training algorithm itself.

\textbf{Agentic Sandbox Environments.}
Multi-turn agent training requires sandboxed environments that provide isolation, reproducibility, and security at scale.
Existing platforms~\citep{wang2024openhands, jimenez2024swebench, jain2025r2egym, yang2024sweagent} have established primary protocols, but they deeply rely on Docker for agent execution.
Docker assumes daemon access and root-equivalent privileges, which are often unavailable on shared Slurm-managed HPC clusters. 
As a result, practitioners often face a trade-off between maintaining separate infrastructure for evaluation and deployment, or incurring the operational complexity of privileged container runtimes on restricted systems.
\prorlagent addresses this limitation by building its sandbox infrastructure on Singularity, enabling rootless execution and native Slurm integration for large-scale agent training on HPC systems.
